\section{Background and related work}
\paragraph{Small models and their deployment.} Surveys catalogue architectures, training and compression techniques for SLMs~\citep{nguyen2024slmsurvey}. Model vendors report
results for their own models on public benchmarks, for example the Qwen2.5-Coder family~\citep{hui2024qwen25coder}; such numbers are measured on server hardware, with
vendor-chosen prompts and sampling. Our focus is the complementary question of what a given quantized checkpoint delivers on consumer hardware.

\paragraph{Quantization and the GGUF format.} Post-training quantization stores weights with 3--8 bits and typically loses little accuracy at 4 bits~\citep{frantar2023gptq,dettmers2023case}; 4-bit precision was found to be close to optimal in total model bits versus zero-shot accuracy~\citep{dettmers2023case}.
llama.cpp~\citep{llamacpp} stores models in GGUF, a single file with weights, tokenizer and metadata, and offers block-wise ``k-quant'' formats such as \code{Q4\_K\_M}, which we use for every model so that
differences are not due to precision.

\paragraph{Inference speed and speculative decoding.} Autoregressive decoding at batch size one is bound by memory bandwidth, because every weight is read for each generated token; the
roofline model~\citep{williams2009roofline} gives the corresponding bound. Speculative decoding drafts several tokens and verifies them in a single pass of the target model, leaving the output distribution
unchanged~\citep{leviathan2023speculative}. Multi-token prediction trains additional output heads to predict several future tokens~\citep{gloeckle2024mtp}; such heads can serve as the draft model.
We measure this on one MTP-enabled checkpoint.

\paragraph{Execution-based evaluation of code.} Functional correctness against hidden tests is the standard for code generation~\citep{chen2021codex}, and insufficient tests are known to
overestimate model quality and mis-rank models~\citep{liu2023evalplus}. We follow this practice, extend it to shell, SQL and \code{kubectl} tasks, and spend effort on validating the graders themselves.

\paragraph{Self-repair and retry loops.} Models can revise their output given feedback~\citep{madaan2023selfrefine,chen2024selfdebug}, but gains from \emph{intrinsic} self-correction are doubtful and often vanish once the
extra compute is accounted for~\citep{huang2024selfcorrect,olausson2024selfrepair}. In particular, \citet{olausson2024selfrepair} found that gains from self-repair are small and uneven once the cost of repair is counted, and attributed this to
the quality of the feedback. Our setting differs in that the feedback comes from an external verifier, and we compare it directly with a no-feedback resampling control on the same first attempts.

\paragraph{Sampling settings and statistical reporting.} Sampling temperature between 0 and 1 had no significant effect on multiple-choice problem solving across nine models~\citep{renze2024temperature}; dynamic truncation such as min-$p$ is
reported to help mainly at high temperature~\citep{nguyen2024minp}. We test vendor-recommended settings against ours directly. On reporting, \citet{miller2024errorbars} argue that evaluations are experiments and should carry error bars, clustered by question and
paired across models; we adopt that approach (bootstrap~\citep{efron1993bootstrap}, Holm correction~\citep{holm1979}, McNemar's test for paired outcomes~\citep{mcnemar1947}).
